% vim: set ts=4 sw=4 tw=80 noexpandtab:

\documentclass{42-en}

%******************************************************************************%
%                                                                              %
%                                   Prologue                                   %
%                                                                              %
%******************************************************************************%
\usepackage[
    type={CC},
    modifier={by-nc-sa},
    version={4.0},
]{doclicense}
\usemintedstyle{monokai}


%****************************************************************%
%                  Re/definition of commands                     %
%****************************************************************%

\newcommand{\ailogo}[1]{\def \@ailogo {#1}}\ailogo{assets/42ai_logo.pdf}

%%  Redefine \maketitle
\makeatletter
\def \maketitle {
  \begin{titlepage}
    \begin{center}
	%\begin{figure}[t]
	  %\includegraphics[height=8cm]{\@ailogo}
	  \includegraphics[height=8cm]{assets/42ai_logo.pdf}
	%\end{figure}
      \vskip 5em
      {\huge \@title}
      \vskip 2em
      {\LARGE \@subtitle}
      \vskip 4em
    \end{center}
    %\begin{center}
	  %\@author
    %\end{center}
	%\vskip 5em
  \vfill
  \begin{center}
    \emph{\summarytitle : \@summary}
  \end{center}
  \vspace{2cm}
  %\vskip 5em
  %\doclicenseThis
  \end{titlepage}
}
\makeatother

\makeatletter
\def \makeheaderfilesforbidden
{
  \noindent
  \begin{tabularx}{\textwidth}{|X X  X X|}
    \hline
  \multicolumn{1}{|>{\raggedright}m{1cm}|}
  {\vskip 2mm \includegraphics[height=1cm]{assets/42ai_logo.pdf}} &
  \multicolumn{2}{>{\centering}m{12cm}}{\small Exercise : \@exnumber } &
  \multicolumn{1}{ >{\raggedleft}p{1.5cm}|}
%%              {\scriptsize points : \@exscore} \\ \hline
              {} \\ \hline

  \multicolumn{4}{|>{\centering}m{15cm}|}
              {\small \@extitle} \\ \hline

  \multicolumn{4}{|>{\raggedright}m{15cm}|}
              {\small Turn-in directory : \ttfamily
                $ex\@exnumber/$ }
              \\ \hline
  \multicolumn{4}{|>{\raggedright}m{15cm}|}
              {\small Files to turn in : \ttfamily \@exfiles }
              \\ \hline

  \multicolumn{4}{|>{\raggedright}m{15cm}|}
              {\small Forbidden functions : \ttfamily \@exforbidden }
              \\ \hline

%%  \multicolumn{4}{|>{\raggedright}m{15cm}|}
%%              {\small Remarks : \ttfamily \@exnotes }
%%              \\ \hline
\end{tabularx}
%% \exnotes
\exrules
\exmake
\exauthorize{None}
\exforbidden{None}
\extitle{}
\exnumber{}
}
\makeatother


\begin{document}

% =============================================================================%
%                     =====================================                    %

\title{Python \& ML - Module 00}
\subtitle{Basic stuff - Eleven Commandments}
\author{
	Maxime Choulika (cmaxime), Pierre Peigné (ppeigne), Matthieu David (mdavid)
}

\summary
{
	The goal of the module is to get started with the Python language.
}

\maketitle
%******************************************************************************%
%                                                                              %
%                        Common Instructions                                   %
%                          for Python Projects                                 %
%                                                                              %
%******************************************************************************%

\chapter{Common Instructions}
\begin{itemize}
  \item The version of Python recommended to use is 3.7, you can
  check the version of Python with the following command: \texttt{python -V}
  
  \item The norm: during this piscine, it is recommended to follow the
  \href{https://www.python.org/dev/peps/pep-0008/}{PEP 8 standards}, though it is not mandatory.
  You can install \href{https://pypi.org/project/pycodestyle}{pycodestyle} which
  is a tool to check your Python code.
  
  \item The function \texttt{eval} is never allowed.
  
  \item The exercises are ordered from the easiest to the hardest.
  
  \item Your exercises are going to be evaluated by someone else,
  so make sure that your variable names and function names are appropriate and civil. 
  
  \item Your manual is the internet.
  
  \item You can also ask questions in the \texttt{\#bootcamps} channel in the \href{https://42-ai.slack.com}{42AI}
  or \href{42born2code.slack.com}{42born2code}.
  
  \item If you find any issue or mistakes in the subject please create an issue on \href{https://github.com/42-AI/bootcamp_python/issues}{42AI repository on Github}.  
  
  \item We encourage you to create test programs for your
  project even though this work \textbf{won't have to be
  submitted and won't be graded}. It will give you a chance
  to easily test your work and your peers’ work. You will find
  those tests especially useful during your defence. Indeed,
  during defence, you are free to use your tests and/or the
  tests of the peer you are evaluating.
  
  \item Submit your work to your assigned git repository. Only the work in the
  git repository will be graded. If Deepthought is assigned to grade your
  work, it will be run after your peer-evaluations.
  If an error happens in any section of your work during Deepthought's grading,
  the evaluation will stop.
\end{itemize}
\newpage
\tableofcontents
\startexercices

%                     =====================================                    %
% =============================================================================%


%******************************************************************************%
%                                                                              %
%                                   Exercises                                  %
%                                                                              %
%******************************************************************************%

% ============================================== %
% ===========================(start ex 00)       %
\chapter{Exercise 00}
\extitle{\$PATH}
\turnindir{ex00}
\exnumber{00}
\exfiles{answers.txt, requirements.txt}
\exforbidden{None}
\makeheaderfilesforbidden

\emph{The first thing you need to do is install Python.}
\\\\
Most modern Unix-based system have a \texttt{python} interpreter installed by default,
but its version might be lower/higher than the one used for these modules.
It is also possible that the default \texttt{python} command uses a version 2.x (for legacy reasons).
This is obviously very confusing for new developper.
\\

\begin{42console}
$> python -V
$> python3 -V
\end{42console}
\hfill \break
To deal with those version issues we will use \texttt{conda}. This program allow you to manage your Python packages and different working environments. 
\\\\
\emph{Note: the actual requirement is to use a Python 3.7.X version. You are free to use a different program/utilities to achieve this goal. At your own risk.}
\pagebreak

% ===========================(Conda Manual)       %
\section*{Conda manual installation}
\emph{Go the next section for an automated installation.}
\\\\
We recommend the following path for your \texttt{conda} folder.
\begin{42console}
$> MYPATH="/goinfre/$USER/miniconda3"
\end{42console}

% ---------------------------------- %
\subsection*{1. Download \& Install conda}

\begin{42console}
# For MAC
$> curl -LO "https://repo.anaconda.com/miniconda/Miniconda3-latest-MacOSX-x86_64.sh"
$> sh Miniconda3-latest-MacOSX-x86_64.sh -b -p $MYPATH

# For Linux
$> curl -LO "https://repo.anaconda.com/miniconda/Miniconda3-latest-Linux-x86_64.sh"
$> sh Miniconda3-latest-Linux-x86_64.sh -b -p $MYPATH
\end{42console}



% ---------------------------------- %
\subsection*{2. Initial configuration of conda}
\begin{42console}
# For zsh
$> $MYPATH/bin/conda init zsh
$> $MYPATH/bin/conda config --set auto_activate_base false
$> source ~/.zshrc

# For bash
$> $MYPATH/bin/conda init bash
$> $MYPATH/bin/conda config --set auto_activate_base false
$> source ~/.bash_profile
\end{42console}



% ---------------------------------- %
\subsection*{3. Create an environment for 42AI !}
\begin{42console}
$> conda create --name 42AI-$USER python=3.7 jupyter pandas pycodestyle numpy
\end{42console}



% ---------------------------------- %
\subsection*{4. Check your 42AI Python environment}
\begin{42console}
$> conda info --envs
$> conda activate 42AI-$USER
$> which python
$> python -V
$> python -c "print('Hello World!')"
\end{42console}



% ---------------------------------- %
\subsection*{Help !}
\begin{itemize}
	\item \textbf{I have lost my miniconda3 folder !} Repeat step 1 and 3.
	\item \textbf{I have lost my home directory !} Repeat step 2.
\end{itemize}
\pagebreak



% ===========================(Conda Auto)       %
\section*{Conda automated installation}

Copy the following script into a file, launch it and follow the instruction. It downloads miniconda, installs it in a \texttt{/goinfre} subfolder and creates a \texttt{python} environment in \texttt{conda}.


\begin{42console}
#!/bin/bash

function which_dl {
	# If operating system name contains Darwnin: MacOS. Else Linux
	if uname -s | grep -iqF Darwin; then
		echo "Miniconda3-latest-MacOSX-x86_64.sh"
	else
		echo "Miniconda3-latest-Linux-x86_64.sh"
	fi
}

function which_shell {
	# if $SHELL contains zsh, zsh. Else Bash
	if echo $SHELL | grep -iqF zsh; then
		echo "zsh"
	else
		echo "bash"
	fi
}

function when_conda_exist {
	# check and install 42AI environement
	printf "Checking 42AI-$USER environment: "
	if conda info --envs | grep -iqF 42AI-$USER; then
		printf "\e[33mDONE\e[0m\n"
	else
		printf "\e[31mKO\e[0m\n"
		printf "\e[33mCreating 42AI environnment:\e[0m\n"
		conda update -n base -c defaults conda -y
		conda create --name 42AI-$USER python=3.7 jupyter numpy pandas pycodestyle -y
	fi
}

function set_conda {
	MINICONDA_PATH="/goinfre/$USER/miniconda3"
	CONDA=$MINICONDA_PATH"/bin/conda"
	PYTHON_PATH=$(which python)
	REQUIREMENTS="jupyter numpy pandas pycodestyle"
	SCRIPT=$(which_dl)
	MY_SHELL=$(which_shell)
	DL_LINK="https://repo.anaconda.com/miniconda/"$SCRIPT
	DL_LOCATION="/tmp/"

	printf "Checking conda: "
	TEST=$(conda -h 2>/dev/null)
	if [ $? == 0 ] ; then
		printf "\e[32mOK\e[0m\n"
		when_conda_exist
		return
	fi
	printf "\e[31mKO\e[0m\n"
	if [ ! -f $DL_LOCATION$SCRIPT ]; then
		printf "\e[33mDonwloading installer:\e[0m\n"
		cd $DL_LOCATION
		curl -LO $DL_LINK
		cd -
	fi
	printf "\e[33mInstalling conda:\e[0m\n"
	sh $DL_LOCATION$SCRIPT -b -p $MINICONDA_PATH
	printf "\e[33mConda initial setup:\e[0m\n"
	$CONDA init $MY_SHELL
	$CONDA config --set auto_activate_base false

	printf "\e[33mCreating 42AI-$USER environnment:\e[0m\n"
	$CONDA update -n base -c defaults conda -y
	$CONDA create --name 42AI-$USER python=3.7 jupyter numpy pandas pycodestyle -y
	printf "\e[33mLaunch the following command or restart your shell:\e[0m\n"
	if [ $MY_SHELL == "zsh" ]; then
		printf "\tsource ~/.zshrc\n"
	else
		printf "\tsource ~/.bash_profile\n"
	fi
}

set_conda
\end{42console}
\hfill \break
Don't forget to check your 42AI Python environment !
\begin{42console}
conda info --envs
conda activate 42AI-$USER
which python
python -V
python -c "print('Hello World!')"
\end{42console}
\newpage



% =================================(Question) %
\section*{(Finally) getting started}

Now that your setup is ready to run, here's a few questions that need to be solved using \texttt{python}, \texttt{pip} or \texttt{conda}. Save your answers in a file \texttt{answers.txt} (one answer per line and per question), and check them with your peers.\\\\
Find the commands to:
\begin{itemize}
	\item Output a list of installed packages and their versions.
	\item Show the package metadata of \texttt{numpy}.
	\item Remove the package \texttt{numpy}.
	\item (Re)install the package \texttt{numpy}.
	\item Freeze your \texttt{python} packages and their versions in a \texttt{requirements.txt} file you have to turn-in.
\end{itemize}


% ===========================(fin ex 00)         %
\newpage
% ===========================(start ex 01)       %
\chapter{Exercise 01}
\extitle{Rev Alpha}
\turnindir{ex01}
\exnumber{01}
\exfiles{exec.py}
\exforbidden{None}
\makeheaderfilesforbidden


% ================================= %
% \section*{}
Make a program that takes a string as argument, reverses it, swaps its letters case and print the result.
\begin{itemize}
	\item If more than one argument are provided, merge them into a single string with each argument separated by a space character.
	\item If no argument are provided, do nothing or print an usage.
\end{itemize}


% --------------------------------- %
\section*{Examples}

\begin{42console}
$> python3 exec.py 'Hello World!' | cat -e
!DLROw OLLEh$
$>
$> python3 exec.py 'Hello' 'my Friend' | cat -e
DNEIRf YM OLLEh$
$>
$> python3 exec.py
$>
\end{42console}

% ===========================(fin ex 01)         %
\newpage
% ===========================(start ex 02)       %
\chapter{Exercise 02}
\extitle{The Odd, the Even and the Zero}
\turnindir{ex02}
\exnumber{02}
\exfiles{whois.py}
\exforbidden{None}
\makeheaderfilesforbidden


% ================================= %
% \section*{}
Make a program that takes a number as argument, checks whether it is odd, even or zero, and print the result.
\begin{itemize}
	\item If more than one argument are provided or if the argument is not an integer, print an error message.
	\item If no argument are provided, do nothing or print an usage.
\end{itemize}


% --------------------------------- %
\section*{Examples}
\begin{42console}
$> python3 whois.py 12
I'm Even.
$>
$> python3 whois.py 3
I'm Odd.
$>
$> python3 whois.py
$>
$> python3 whois.py 0
I'm Zero.
$>
$> python3 whois.py Hello
AssertionError: argument is not an integer
$>
$> python3 whois.py 12 3
AssertionError: more than one argument are provided
$>
\end{42console}

\hint{
	No bonus point to be gained from a complex error management system. Keep it simple.
}

% ===========================(fin ex 02)         %
\newpage
% ===========================(start ex 03)       %
\chapter{Exercise 03}
\extitle{Elementary}
\turnindir{ex03}
\exnumber{03}
\exfiles{operations.py}
\exforbidden{None}
\makeheaderfilesforbidden


% ================================= %
% \section*{}
Write a program that takes two integers A and B as arguments and prints the result of the following operations:
\begin{42console}
Sum:         A+B
Difference: A-B
Product:     A*B
Quotient:    A/B
Remainder:   A%B
\end{42console}

\begin{itemize}
	\item If more or less than two argument are provided or if either of the argument is not an integer, print an error message.
	\item If no argument are provided, do nothing or print an usage.
	\item If an operation is impossible, print an error message instead of a numerical result.
\end{itemize}

% ================================= %
\section*{Examples}
\begin{42console}
$> python3 operations.py 10 3
Sum:         13
Difference:  7
Product:     30
Quotient:    3.3333...
Remainder:   1
$>
$> python3 operations.py 42 10
Sum:         52
Difference:  32
Product:     420
Quotient:    4.2
Remainder:   2
$>
$> python3 operations.py 1 0
Sum:         1
Difference:  1
Product:     0
Quotient:    ERROR (division by zero)
Remainder:   ERROR (modulo by zero)
$>
$> python3 operations.py
Usage: python operations.py <number1> <number2>
Example:
	python operations.py 10 3
$>
$> python3 operations.py 12 10 5
AssertionError: too many arguments
$>
$> python3 operations.py "one" "two"
AssertionError: only integers
$>
\end{42console}

\hint{
	No bonus point to be gained from handling decimal point or scientific notation. Keep it simple.
}

% ===========================(fin ex 03)         %
\newpage
% ===========================(start ex 04)       %
\chapter{Exercise 04}
\extitle{Functional file}
\turnindir{ex04}
\exnumber{04}
\exfiles{count.py}
\exforbidden{None}
\makeheaderfilesforbidden


% ================================= %
\section*{Part 1. text\_analyzer}
Create a function called \texttt{text\_analyzer} that takes a single string argument
and displays the sums of its upper-case characters, lower-case characters, punctuation characters and spaces.
\begin{itemize}
	\item If \texttt{None} or nothing is provided, the user is prompted to provide a string. 
	\item If the argument is not a string, print an error message.
	\item This function must have a \texttt{docstring} explaning its behavior.
\end{itemize}

Test your function with the \texttt{python} console

% --------------------------------- %
\subsection*{Examples}
\begin{42console}
$> python3
>>> from count import text_analyzer
>>> text_analyzer("Python 2.0, released 2000, introduced
features like List comprehensions and a garbage collection
system capable of collecting reference cycles.")
The text contains 143 character(s):
- 2 upper letter(s)
- 113 lower letter(s)
- 4 punctuation mark(s)
- 18 space(s)
>>> text_analyzer("Python is an interpreted, high-level,
general-purpose programming language. Created by Guido van
Rossum and first released in 1991, Python's design philosophy
emphasizes code readability with its notable use of significant
whitespace.")
The text contains 234 character(s):
- 5 upper letter(s)
- 187 lower letter(s)
- 8 punctuation mark(s)
- 30 space(s)
>>> text_analyzer()
What is the text to analyze?
>> Hello World!
The text contains 8 character(s):
- 2 upper letter(s)
- 8 lower letter(s)
- 1 punctuation mark(s)
- 1 space(s)
>>> text_analyzer(42)
AssertionError: argument is not a string
>>> print(text_analyzer.__doc__)

    This function counts the number of upper characters, lower characters,
    punctuation and spaces in a given text.
\end{42console}

% ================================= %
\section*{Part 2. \_\_name\_\_==\_\_main\_\_}


In the previous part, you wrote a function that can be used in the console or in another file when  imported.
Without changing this behavior, update your file so it can also be launched as a standalone program.

\begin{itemize}
	\item If more than one argument is provided to the program, print an error message. 
	\item Otherwise, use the \texttt{text\_analyzer} function.
\end{itemize}

% --------------------------------- %
\subsection*{Examples}
\begin{42console}
$> python3 count.py 'Hello World!'
The text contains 8 character(s):
- 2 upper letter(s)
- 8 lower letter(s)
- 1 punctuation mark(s)
- 1 space(s)
$> python3
>>> from count import text_analyzer
>>> text_analyzer("Hello World!")
The text contains 8 character(s):
- 2 upper letter(s)
- 8 lower letter(s)
- 1 punctuation mark(s)
- 1 space(s)
\end{42console}

% ===========================(fin ex 04)         %
\newpage
% ===========================(start ex 05)       %
\chapter{Exercise 05}
\extitle{A recipe}
\turnindir{ex05}
\exnumber{05}
\exfiles{recipe.py}
\exforbidden{None}
\makeheaderfilesforbidden

% ================================= %
\section*{Part 1: Nested Dictionaries}

Create a dictionary called \texttt{cookbook}. You will use this \texttt{cookbook} to store recipe.
\\\\
A recipe is a \textbf{dictionary} that stores (at least) 3 couples key-value:
\begin{itemize}
	\item ''ingredients": a \textbf{list of string} representing the list of ingredients
	\item "meal": a \textbf{string} representing the type of meal
	\item "prep\_time": a \textbf{non-negative integer} representing a time in minutes
\end{itemize}

In the \texttt{cookbook}, the  \textbf{key} to a recipe is the recipe name.
\\
Initialize your \texttt{cookbook} with 3 recipes:
\begin{itemize}
	\item The Sandwich's ingredients are \textit{ham}, \textit{bread}, \textit{cheese} and \textit{tomatoes}.
It is a \textit{lunch} and it takes $10$ minutes of preparation.
	\item The Cake's ingredients are \textit{flour}, \textit{sugar} and \textit{eggs}.
It is a \textit{dessert} and it takes $60$ minutes of preparation.
	\item The Salad's ingredients are \textit{avocado}, \textit{arugula}, \textit{tomatoes} and \textit{spinach}.
It is a \textit{lunch} and it takes $15$ minutes of preparation.
\end{itemize}

% ================================= %
\section*{Part 2: A series of Helpful Functions}

Create a series of useful functions to handle your \texttt{cookbook}:

\begin{enumerate}
	\item A function that print all recipe names.
	\item A function that takes a recipe name and print its details.
	\item A function that takes a recipe name and delete it.
	\item A function that add a recipe from user input. You will need a name, a list of ingredient, a meal type and a preparation time.
\end{enumerate}
\subsection*{input example}
\begin{42console}
>>> Enter a name:
chips
>>> Enter ingredients:
potatoes
oil
salt

>>> Enter a meal type:
lunch
>>> Enter a preparation time:
15
\end{42console}

% ================================= %
\section*{Part 3: A command line executable !}

Create a program that use your \texttt{cookbook} and your functions.\\

The program will prompt the user to make a choice between printing the cookbook content, printing one recipe, adding a recipe, deleting a recipe or quitting the cookbook.\\

Your program will continue to ask for prompt until the user decide to quit it. The program cannot crash if a wrong value is entered: you must handle the error and ask for another prompt.

\begin{42console}
$> python3 recipe.py
Welcome to the Python Cookbook !
List of available option:
   1: Add a recipe
   2: Delete a recipe
   3: Print a recipe
   4: Print the cookbook
   5: Quit

Please select an option:
>> 3

Please enter a recipe name to get its details:
>> cake

Recipe for cake:
   Ingredients list: ['flour', 'sugar', 'eggs']
   To be eaten for dessert.
   Takes 60 minutes of cooking.

Please select an option:
>> Hello

Sorry, this option does not exist.
List of available option:
   1: Add a recipe
   2: Delete a recipe
   3: Print a recipe
   4: Print the cookbook
   5: Quit

Please select an option:
>> 5

Cookbook closed. Goodbye !
$>
\end{42console}


% ===========================(fin ex 05)         %
\newpage
% ===========================(start ex 06)       %
\chapter{Exercise 06}
\extitle{Shorter, faster, pythonest}
\turnindir{ex06}
\exnumber{06}
\exfiles{filterwords.py}
\exforbidden{filter}
\makeheaderfilesforbidden

Make a program that takes a string S and an integer N as argument and print the list of words in S that contains more than N non-punctuation characters.

\begin{itemize}
	\item Words are separated from each other by space characters
	\item Punctuation symbols must be removed from the printed list: they are neither part of a word nor a separator
	\item The program must contains at least one \textbf{list comprehension} expression. 
\end{itemize}

If the number of argument is different from 2, or if the type of any argument is wrong, the program prints an error message.

% ================================= %
\section*{Examples}
% --------------------------------- %

\begin{42console}
$> python3 filterwords.py 'Hello, my friend' 3
['Hello', 'friend']
$> python3 filterwords.py 'Hello, my friend' 10
[]
$>  python3 filterwords.py 'A robot must protect its own existence as long as such protection does not conflict with the First or Second Law' 6
['protect', 'existence', 'protection', 'conflict']
$> python3 filterwords.py Hello World
ERROR
$> python3 filterwords.py 3 'Hello, my friend'
ERROR
$> python3 filterwords.py
ERROR
\end{42console}

% ===========================(fin ex 06)         %
\newpage
% ===========================(start ex 07)       %

\chapter{Exercise 07}
\extitle{The right format}
\turnindir{ex07}
\exnumber{07}
\exfiles{kata00.py, kata01.py, kata02.py, kata03.py, kata04.py}
\exforbidden{None}
\makeheaderfilesforbidden

Let's get familiar with the useful concept of \textbf{string formatting} through a kata series.\\

Each exercice will provide you with a \texttt{kata} variable. This variable can be modified to a certain extent: your program must react accordingly.

% ================================= %
\section*{kata00}
The \texttt{kata} variable is always a tuple and can only be filled with integer.
\begin{42console}
# Put this at the top of your kata00.py file
kata = (19,42,21)
\end{42console}

Write a program that display this variable content according to the format shown below:

\begin{42console}
$> python3 kata00.py
The 3 numbers are: 19, 42, 21
$>
\end{42console}

% ================================= %
\section*{kata01}

The \texttt{kata} variable is always a dictionary and can only be filled with strings.

\begin{42console}
# Put this at the top of your kata01.py file
kata = {
    'Python': 'Guido van Rossum',
    'Ruby': 'Yukihiro Matsumoto',
    'PHP': 'Rasmus Lerdorf',
    }
\end{42console}

Write a program that display this variable content according to the format shown below:

\begin{42console}
$> python3 kata01.py
Python was created by Guido van Rossum
Ruby was created by Yukihiro Matsumoto
PHP was created by Rasmus Lerdorf
$>
\end{42console}

% ================================= %
\section*{kata02}

The \texttt{kata} variable is always a tuple that contains 5 non-negative integers. The first integer contains up to 4 digits, the rest up to 2 digits.

\begin{42console} 
# Put this at the top of your kata02.py file
kata = (2019, 9, 25, 3, 30)
\end{42console}

Write a program that display this variable content according to the format shown below:

\begin{42console}
$> python3 kata02.py | cat -e
09/25/2019 03:30$
$> python3 kata02.py | wc -c
17
$>
\end{42console}

% ================================= %
\section*{kata03}

The \texttt{kata} variable is always a string whose length is not higher than 42.

\begin{42console}
# Put this at the top of your kata03.py file
kata = "The right format"
\end{42console}

Write a program that display this variable content according to the format shown below:

\begin{42console}
$> python3 kata03.py | cat -e
--------------------------The right format%
$> python3 kata03.py | wc -c
42
$>
\end{42console}

% ================================= %
\section*{kata04}

The \texttt{kata} variable is always a tuple that contains, in the following order:
\begin{itemize}
	\item 2 non-negative integer containing up to 2 digits
	\item 1 decimal
	\item 1 integer
	\item 1 decimal
\end{itemize}

\begin{42console}
# Put this at the top of your kata04.py file
kata = (0, 4, 132.42222, 10000, 12345.67)
\end{42console}

Write a program that display this variable content according to the format shown below:

\begin{42console}
$> python3 kata04.py
module_00, ex_04 : 132.42, 1.00e+04, 1.23e+04
$> python3 kata04.py | cut -c 10,18
,:
\end{42console}

% ===========================(fin ex 07)         %
\newpage
% ============================================== %
% ===========================(start ex 08)       %
\chapter{Exercise 08}
\extitle{Loading bar!}
\turnindir{ex08}
\exnumber{08}
\exfiles{loading.py}
\exforbidden{\texttt{tqdm} or any library for automatic loading bar}
\makeheaderfilesforbidden


You are about to discover the \texttt{yield} operator!

So let's create a function called \texttt{ft\_progress(lst)}.

The function will display the progress of a \texttt{for} loop.


% ================================= %
\section*{Examples}
% --------------------------------- %
\begin{42console}% SERAIT MIEUX AVEC UN ENVIRONNEMENT PYTHON
	listy = range(1000)
	ret = 0
	for elem in ft_progress(listy):
	    ret += (elem + 3) % 5
	    sleep(0.01)
	print()
	print(ret)
\end{42console}

\begin{42console}
	$> python loading.py
	ETA: 8.67s [ 23%][=====>                  ] 233/1000 | elapsed time 2.33s
	...
	2000
\end{42console}

\begin{42console}% SERAIT MIEUX AVEC UN ENVIRONNEMENT PYTHON
	listy = range(3333)
	ret = 0
	for elem in ft_progress(listy):
		ret += elem
		sleep(0.005)
	print()
	print(ret)
\end{42console}

\begin{42console}
	$> python loading.py
	ETA: 14.67s [  9%][=>                      ] 327/3333 | elapsed time 1.33s
	...
	5552778
\end{42console}

\info{We advise you to go take a look at the wonderful tqdm library, it will come in handy in many situations}

% ===========================(fin ex 08)         %
% ============================================== %
\newpage

% ============================================== %
% ===========================(start ex 09)       %
\chapter{Exercise 09}
\extitle{Secret number}
\turnindir{ex09}
\exnumber{09}
\exfiles{guess.py}
\exforbidden{None}
\makeheaderfilesforbidden

You have to make a program that will be an interactive guessing game.
It will ask the user to guess a number between $1$ and $99$.
The program will tell the user if their input is too high or too low.
The game ends when the user finds out the secret number or types \texttt{exit}.
You will import the \texttt{random} module with the \texttt{randint} function to get a random number.
You have to count the number of trials and print that number when the user wins.

% ================================= %
\section*{Examples}
% --------------------------------- %

\begin{42console}
	$> python guess.py
	This is an interactive guessing game!
	You have to enter a number between 1 and 99 to find out the secret number.
	Type 'exit' to end the game.
	Good luck!

	What's your guess between 1 and 99?
	>> 54
	Too high!
	What's your guess between 1 and 99?
	>> 34
	Too low!
	What's your guess between 1 and 99?
	>> 45
	Too high!
	What's your guess between 1 and 99?
	>> A
	That's not a number.
	What's your guess between 1 and 99?
	>> 43
	Congratulations, you've got it!
	You won in 5 attempts!
\end{42console}

If the user discovers the secret number on the first try, tell them.
If the secret number is 42, make a reference to Douglas Adams.

\begin{42console}
	$> python guess.py
	This is an interactive guessing game!
	You have to enter a number between 1 and 99 to find out the secret number.
	Type 'exit' to end the game.
	Good luck!

	What's your guess between 1 and 99?
	>> 42
	The answer to the ultimate question of life, the universe and everything is 42.
	Congratulations! You got it on your first try!
\end{42console}


Other example:

\begin{42console}
	$> python guess.py
	This is an interactive guessing game!
	You have to enter a number between 1 and 99 to find out the secret number.
	Type 'exit' to end the game.
	Good luck!

	What's your guess between 1 and 99?
	>> exit
	Goodbye!
\end{42console}

% ===========================(fin ex 09)         %
% ============================================== %

\newpage

% ===========================(start ex 10)       %
\chapter{Exercise 10}
\extitle{S.O.S}
\turnindir{ex10}
\exnumber{10}
\exfiles{sos.py}
\exforbidden{None}
\makeheaderfilesforbidden

Make a program that takes a string as argument and encode it into Morse code.

\begin{itemize}
	\item The program supports space and alphanumeric characters
	\item An alphanumeric character is represented by dots \texttt{.} and dashes \texttt{-}: 
	\item A space character is represented by a slash \texttt{/}
	\item Complete morse characters are separated by a single space\\
\end{itemize}

If more than one argument are provided, merge them into a single string with each argument separated by a space character.\\

If no argument is provided, do nothing or print an usage.

% ================================= %
\section*{Examples}
\begin{42console}
	$> python3 sos.py "SOS"
	... --- ...
	$> python3 sos.py
	$> python3 sos.py "HELLO / WORLD"
	ERROR
	$> python3 sos.py "96 BOULEVARD" "Bessiere"
	----. -.... / -... --- ..- .-.. . ...- .- .-. -.. / -... . ... ... .. . .-. .
\end{42console}


\hint{
	\url{https://morsecode.world/international/morse2.html}
}
% ===========================(fin ex 10)         %
% ============================================== %

\newpage
% ================================= %
\section*{Contact}
% --------------------------------- %
You can contact 42AI association by email: contact@42ai.fr\\
You can join the association on \href{https://join.slack.com/t/42-ai/shared_invite/zt-ebccw5r7-YPkDM6xOiYRPjqJXkrKgcA}{42AI slack}
and/or apply to \href{https://forms.gle/VAFuREWaLmaqZw2D8}{one of the association teams}.

% ================================= %
\section*{Acknowledgements}
% --------------------------------- %
The modules Python \& ML is the result of a collective work, we would like to thanks:
\begin{itemize}
  \item Maxime Choulika (cmaxime),
  \item Pierre Peigné (ppeigne, pierre@42ai.fr),
  \item Matthieu David (mdavid, matthieu@42ai.fr),
  \item Quentin Feuillade--Montixi (qfeuilla, quentin@42ai.fr)
\end{itemize}
who supervised the creation, the enhancement and this present transcription.

\begin{itemize}
    \item Louis Develle (ldevelle, louis@42ai.fr)
    \item Augustin Lopez (aulopez)
    \item Luc Lenotre (llenotre)
    \item Owen Roberts (oroberts)
    \item Thomas Flahault (thflahau)
    \item Amric Trudel (amric@42ai.fr)
    \item Baptiste Lefeuvre (blefeuvr@student.42.fr)
    \item Mathilde Boivin (mboivin@student.42.fr)
    \item Tristan Duquesne (tduquesn@student.42.fr)
\end{itemize}
for your investment for the creation and development of these modules.

\begin{itemize}
    \item Barthélémy Leveque (bleveque@student.42.fr)
    \item Remy Oster (roster@student.42.fr)
    \item Quentin Bragard (qbragard@student.42.fr)
    \item Marie Dufourq (madufour@student.42.fr)
    \item Adrien Vardon (advardon@student.42.fr)
\end{itemize}
who betatest the first version of the modules of Machine Learning.
\vfill
\doclicenseThis

\end{document}
